\section{Posterior Schr\"odinger Bridges}
\label{sec:method}

Posterior Schr\"odinger Bridges (PSB) models each task transition as a
Schr\"odinger bridge from a carried posterior approximation to its update
under the new-task likelihood. At boundary $t$, the drift fitted at the
preceding boundary guides an annealed traversal to a terminal population.
The terminal networks define the next carried approximation. PSB then fits
a new bridge between the carried approximation and the terminal population.
Its drift guides boundary $t+1$. Thus, every boundary has its own bridge,
while the learned drifts carry information forward across tasks.

\subsection{Posterior endpoints}
\label{sec:representation}

After task $t$, we retain $K$ network parameters $\theta_t^{(k)}$ as the
centres of a Gaussian mixture,
\begin{equation}
    \hat q_t(\theta)=\sum_{k=1}^{K}w_t^{(k)}
    \mathcal{N}\!\left(\theta\mid\theta_t^{(k)},A_t^{-1}\right),
    \label{eq:mixture}
\end{equation}
where the nonnegative weights sum to one and $A_t$ is a shared diagonal
precision. The mixture can retain several parameter solutions, while its
precision limits variation in directions important to earlier tasks. For
task $1$, we train $K$ networks independently from He-uniform
initializations~\cite{he2015delving} and use their final parameters as the
mixture centres with uniform weights. The trained networks define an
approximation to $q_1$. Starting parameters at later boundaries are drawn
from the full mixture in (\ref{eq:mixture}), including its Gaussian spread.

We update the precision using diagonal empirical Fisher information,
\begin{equation}
    A_t=\alpha I+\sum_{j=1}^{t}\gamma^{t-j}N_jF_j,
    \label{eq:precision}
\end{equation}
where $N_j=|\mathcal{D}_j|$, $\gamma\in(0,1]$ discounts older tasks, and
$\alpha>0$ keeps the precision invertible. The estimate for task $j$ is
\begin{equation}
    F_j=\frac{1}{N_j}\sum_{(x,y)\in\mathcal{D}_j}\sum_{k=1}^{K}w_j^{(k)}
    \left(\nabla_\theta\log p(y\mid x,\theta_j^{(k)})\right)^{\odot2}.
    \label{eq:fisher}
\end{equation}
Thus, $N_jF_j$ approximates the diagonal curvature of the full task
likelihood. When $\gamma=1$, the estimates accumulate without discounting.

Given $\hat q_t$, the incoming task defines the updated target
\begin{equation}
    \tilde q_{t+1}(\theta)\propto
    p(\mathcal{D}_{t+1}\mid\theta)\hat q_t(\theta).
    \label{eq:target}
\end{equation}
This applies Bayes' rule to the approximation carried from previous tasks.
Its likelihood is the product over the $N_{t+1}$ examples, so its log is
the sum of their log-likelihoods. To approach the target gradually, define
\begin{equation}
    \pi_{t,\beta}(\theta)\propto
    p(\mathcal{D}_{t+1}\mid\theta)^\beta\hat q_t(\theta),
    \qquad \beta\in[0,1].
    \label{eq:tempered}
\end{equation}
Here $\pi_{t,0}=\hat q_t$ and $\pi_{t,1}=\tilde q_{t+1}$. Writing
$r_k(\theta)$ for the normalized contribution of component $k$ in
(\ref{eq:mixture}), the score is
\begin{equation}
\begin{split}
    \nabla_\theta\log\pi_{t,\beta}(\theta)
    &=\beta\nabla_\theta\log p(\mathcal{D}_{t+1}\mid\theta)\\
    &\quad-A_t\left(\theta-
    \sum_{k=1}^{K}r_k(\theta)\theta_t^{(k)}\right).
\end{split}
    \label{eq:pathscore}
\end{equation}
The two terms measure support from the incoming task and from the carried
mixture. For one component, maximizing (\ref{eq:target}) is equivalent to
minimizing
\begin{equation}
\begin{split}
    \bar\ell_{t+1}(\theta)
    &+\frac{1}{2}\sum_{j=1}^{t}\gamma^{t-j}
    \frac{N_j}{N_{t+1}}\|\theta-\theta_t\|_{F_j}^{2}\\
    &+\frac{\alpha}{2N_{t+1}}\|\theta-\theta_t\|^2,
\end{split}
    \label{eq:ewc}
\end{equation}
where $\bar\ell_{t+1}=-N_{t+1}^{-1}\log
p(\mathcal{D}_{t+1}\mid\theta)$ and $\|u\|_M^2=u^\top Mu$.
The resulting MAP objective has the form of online
EWC~\cite{kirkpatrick2017ewc,schwarz2018progress}, with an additional
isotropic term. PSB instead carries multiple components and estimates the
updated distribution.

\subsection{Fisher-metric Schr\"odinger bridge}
\label{sec:bridge}

At boundary $t$, the ideal bridge connects $\hat q_t$ to
$\tilde q_{t+1}$. We choose a Brownian reference with initial
distribution $\mathbb{Q}_{t,0}=\hat q_t$ and dynamics
\begin{equation}
    \mathbb{Q}_t:\quad
    d\Theta_s=\sigma A_t^{-1/2}\,dW_s,\qquad s\in[0,1],
    \label{eq:fisherreference}
\end{equation}
where $\sigma>0$ sets the noise level. The bridge is
\begin{equation}
    \mathbb{P}_t^\star=
    \mathop{\arg\min}_{\substack{\mathbb{P}:\ \mathbb{P}_0=\hat q_t\\
    \mathbb{P}_1=\tilde q_{t+1}}}
    \mathrm{KL}(\mathbb{P}\,\|\,\mathbb{Q}_t).
    \label{eq:psb}
\end{equation}
If $\mathbb{P}$ adds drift $u_s$ to (\ref{eq:fisherreference}), its path
cost is
\begin{equation}
    \mathrm{KL}(\mathbb{P}\,\|\,\mathbb{Q}_t)
    =\frac{1}{2\sigma^2}\mathbb{E}_{\mathbb{P}}
    \int_0^1\|u_s\|_{A_t}^{2}\,ds.
    \label{eq:fishercost}
\end{equation}
The bridge therefore charges more for motion in directions with greater
accumulated Fisher information. In practice, both endpoint distributions
are represented by finite populations, so the fitted bridge approximates
(\ref{eq:psb}). The bridge time $s$ describes its stochastic path. The
temperature $\beta$ in (\ref{eq:tempered}) describes a separate sequence
used to obtain the terminal population.

\subsection{Learning the bridge drift}
\label{sec:drift}

Given starting and terminal populations, we fit the bridge by iterative
Markovian fitting (IMF)~\cite{shi2023dsbm}. We initially pair samples from
the two populations independently. For a paired start $\Theta_0$ and end
$\Theta_1$, the conditional reference bridge at time $s$ is
\begin{equation}
    \Theta_s=(1-s)\Theta_0+s\Theta_1+
    \sigma\sqrt{s(1-s)}A_t^{-1/2}\xi,
    \quad \xi\sim\mathcal{N}(0,I).
    \label{eq:conditionalbridge}
\end{equation}
Let $\Pi$ denote the current pairing. The forward and backward drift
networks are fitted to the conditional bridge velocities,
\begin{align}
    \mathcal{L}_{\mathrm f}(\phi_t)
    &=\mathbb{E}\left\|v^{\mathrm f}_{\phi_t}(\Theta_s,s)
    -\frac{\Theta_1-\Theta_s}{1-s}\right\|_{A_t}^{2},
    \label{eq:lossf}\\
    \mathcal{L}_{\mathrm b}(\psi_t)
    &=\mathbb{E}\left\|v^{\mathrm b}_{\psi_t}(\Theta_s,s)
    -\frac{\Theta_0-\Theta_s}{s}\right\|_{A_t}^{2}.
    \label{eq:lossb}
\end{align}
The expectation uses $(\Theta_0,\Theta_1)\sim\Pi$, bridge samples from
(\ref{eq:conditionalbridge}), and $s$ sampled away from $0$ and $1$.
We then obtain a new pairing by simulating the fitted process from one
endpoint and repeat the regression, alternating forward and backward
directions. Exact IMF projections preserve the endpoint marginals. Finite
populations and learned drifts preserve them only approximately.

The drifts act on neuron tokens so that one predictor can be shared across
units. A fully connected token contains a unit's incoming weights and bias.
A convolutional token contains one output filter and its bias. Let
$z_{i,s}^{(\ell)}$ be token $i$ of layer $\ell$ at $\Theta_s$. Let
$g_{i,s}^{(\ell)}$ and $a_{i,s}^{(\ell)}$ be the matching parts of the
new-task likelihood score and the carried mixture score in
(\ref{eq:pathscore}). The forward drift is assembled from
\begin{equation}
    \big[v^{\mathrm f}_{\phi_t}(\Theta_s,s)\big]_i^{(\ell)}
    =f_{\phi_t}^{(\ell)}\!\left(
    z_{i,s}^{(\ell)},s,g_{i,s}^{(\ell)},a_{i,s}^{(\ell)}\right).
    \label{eq:tokendrift}
\end{equation}
One network $f_{\phi_t}^{(\ell)}$ is shared across the units of layer
$\ell$. This supplies token-level training examples from each network
while keeping the drift parameterization fixed as tasks arrive. The
backward drift uses the same architecture. Conditioning on the two scores
lets the retained networks respond to the likelihood and prior at a new
boundary, where both scores are recomputed.

\subsection{Successive task transitions}
\label{sec:transition}

At boundary $t$, we draw starting parameters from $\hat q_t$ and raise
$\beta$ through $0=\beta_0<\cdots<\beta_M=1$. The forward drift fitted
at boundary $t-1$, denoted $v^{\mathrm f}_{\phi_{t-1}}$, guides each move.
At the first boundary, we set this drift to zero. With
$s_m=\beta_m$ and $\Delta s_m=s_m-s_{m-1}$, a proposal is
\begin{equation}
\begin{split}
    \Theta_m=\Theta_{m-1}
    &+v^{\mathrm f}_{\phi_{t-1}}(\Theta_{m-1},s_{m-1})\Delta s_m\\
    &+\sigma\sqrt{\Delta s_m}A_t^{-1/2}\xi_m.
\end{split}
    \label{eq:euler}
\end{equation}
This use of $s_m=\beta_m$ evaluates the preceding drift during the
proposal. It does not identify the intermediate bridge distributions
with the tempered targets. At boundary $t$, the inputs to
$v^{\mathrm f}_{\phi_{t-1}}$ use the new-task likelihood score and the
prior score of $\hat q_t$.

Let $K_m$ be the Gaussian density of (\ref{eq:euler}) and let $L_m$ be a
normalized backward transition density obtained from the preceding
backward drift. At the first boundary, both densities use the Brownian
reference. Writing
$\varrho_{t,\beta}(\theta)=p(\mathcal{D}_{t+1}\mid\theta)^\beta
\hat q_t(\theta)$, we update each particle's weight by
\begin{equation}
    \omega_m^{(k)}\propto\omega_{m-1}^{(k)}
    \frac{\varrho_{t,\beta_m}(\Theta_m^{(k)})
    L_m(\Theta_{m-1}^{(k)}\mid\Theta_m^{(k)})}
    {\varrho_{t,\beta_{m-1}}(\Theta_{m-1}^{(k)})
    K_m(\Theta_m^{(k)}\mid\Theta_{m-1}^{(k)})}.
    \label{eq:bridgeweight}
\end{equation}
The correction accounts for the difference between the guided proposal
and the tempered target~\cite{delmoral2006smc}. Before each move, we
choose the largest temperature increment whose provisional likelihood
weights have effective size at least $\rho K$, for a chosen
$\rho\in(0,1)$. For weights $\omega^{(k)}$,
this size is $(\sum_k\omega^{(k)})^2/\sum_k(\omega^{(k)})^2$.
If the effective size of the corrected weights falls below $K/2$, we
draw a new population according to their normalized values. We also apply
preconditioned Langevin moves targeting
$\pi_{t,\beta_m}$,
\begin{equation}
    \theta'=\theta+\eta A_t^{-1}
    \nabla_\theta\log\pi_{t,\beta_m}(\theta)
    +\sqrt{2\eta}A_t^{-1/2}\xi,
    \quad\xi\sim\mathcal{N}(0,I),
    \label{eq:mala}
\end{equation}
with acceptance probabilities computed from the full tempered target
and the forward and reverse proposal densities. At $\beta=1$, a final
draw according to the normalized weights gives an equally weighted
terminal population approximating $\tilde q_{t+1}$.

The terminal networks become the centres of $\hat q_{t+1}$, with
\begin{equation}
    A_{t+1}=\alpha I+\gamma(A_t-\alpha I)+N_{t+1}F_{t+1}.
    \label{eq:precisionupdate}
\end{equation}
PSB then fits the bridge at boundary $t$ between samples from
$\hat q_t$ and the terminal population, using the reference based on
$A_t$. We initialize its drift networks from the preceding boundary's
parameters and retain the fitted $v^{\mathrm f}_{\phi_t}$ and
$v^{\mathrm b}_{\psi_t}$ to guide the next transition. After computing
$F_{t+1}$, we discard $\mathcal{D}_{t+1}$. At test time, we use the
weighted predictions of the component centres,
$p(y\mid x)\approx\sum_k w_{t+1}^{(k)}
p(y\mid x,\theta_{t+1}^{(k)})$. The retained state has fixed size in the
number of tasks.
